% BOOK-FILE-ID: f-d44b8bbd86e8
\section{答えの形を見通す}\label{節：答えの形を見通す}

解法を選び,行き詰まったら切り替える.ここまではその練習をしてきた.
ここでは,もう1つの武器を紹介しよう.解く前に,答えの形を見通すことである.
数項を計算する実験が「値」から答えを予想する方法だとすれば,これは「形」から答えを予想する方法である.

\paragraphbf{一般項を並べてみる}
まずは,本編で得た一般項を並べてみよう.
ここでは特性方程式型の$p$は$0,1$以外,階差等比複合型の$p$は0でないとする.
隣接3項型の特性方程式の解は異なる2つの0でない数とし,定数の付加項がある場合はさらに1を解に持たないとする.
\begin{itemize}
  \item 特性方程式型$a_{n+1}=pa_n+q$：\,$a_n=(a_1-\alpha)p^{n-1}+\alpha$
  \item 階差等比複合型$a_{n+1}=pa_n+f(n)$：\,$a_n=\{a_1-g(1)\}p^{n-1}+g(n)$
  \item 隣接3項型$a_{n+2}+pa_{n+1}+qa_n=0$：\,$a_n=A\alpha^{n-1}+B\beta^{n-1}$
  \item 定数の付加項がある隣接3項型：\,$a_n=A\alpha^{n-1}+B\beta^{n-1}+\gamma$
\end{itemize}
じっと眺めると,どれも同じ形をしていないだろうか.
\begin{center}
  (初項で決まる定数)$\times$(付加項のない式の解)\\
  $+$(付加項に合わせて決まる1つの解)
\end{center}
という形である.
例えば特性方程式型なら,\,$p^{n-1}$は付加項$q$を取り去った$b_{n+1}=pb_n$の解であり,\,$\alpha$は元の漸化式を満たす定数の解である.

\paragraphbf{なぜ同じ形になるのか}
これらの漸化式は,どれも項を定数倍して足しただけの形をしている.
このような漸化式では,同じ漸化式を満たす2つの数列$a_n,v_n$の差を取ると,付加項がきれいに消える.
例えば$a_{n+1}=pa_n+f(n)$なら,
\BookMathLayout{\[
  \begin{array}{rrl}
    &a_{n+1}={}&pa_n+f(n)\\
    -\,)&v_{n+1}={}&pv_n+f(n)\\ \hline
    &a_{n+1}-v_{n+1}={}&p(a_n-v_n)
  \end{array}
\]}{\[
  \begin{array}{rl}
    &a_{n+1}=pa_n+f(n)\\
    -\,)&v_{n+1}=pv_n+f(n)\\ \hline
    &a_{n+1}-v_{n+1}\\
    &\quad=p(a_n-v_n)
  \end{array}
\]}
である.つまり,\,$a_n-v_n$は付加項のない漸化式を満たす.
したがって,1つの解$v_n$さえ見つければ,あとは付加項のない漸化式を解けばよいのだ.
本編で不動点や不動関数を引く方法も,この操作にあたる.

次に,付加項のない漸化式\footnote{このような漸化式を\bookterm{homogeneous}{同次}であるという.}を考える.
例えば$a_{n+2}+pa_{n+1}+qa_n=0$の解$u_n,w_n$があれば,定数$C,D$で作った$Cu_n+Dw_n$も
\begin{align*}
  &C(u_{n+2}+pu_{n+1}+qu_n)\\
  &\quad+D(w_{n+2}+pw_{n+1}+qw_n)=0
\end{align*}
より,また解になる.\footnote{この性質を\bookterm{superposition}{重ね合わせの原理}といい,\,$Cu_n+Dw_n$を$u_n,w_n$の\bookterm{combination}{線形結合}という.}
では,いくつの解を組み合わせればよいだろうか.
答えは,自由に決められる初期値の個数である.
2項間の漸化式なら初項$a_1$だけを決めればよいので定数は1つ,3項間なら$a_1,a_2$を決めるので定数は2つになる.
隣接3項型の$a_n=A\alpha^{n-1}+B\beta^{n-1}$で,\,$A,B$を$a_1,a_2$から決めたのはこのためである.\footnote{0でない重解のときは$\alpha^{n-1}$と$(n-1)\alpha^{n-1}$の2つを組み合わせる.0が解の場合は\sectionref*{節：付加項がない場合}を,根が3重,4重と重なる場合は\sectionref*{節：差分で線形漸化式を解く}を参照.}

\begin{bookpoint}
  \textbf{POINT}\par
  項を定数倍して足しただけの漸化式の一般項は,「付加項に合わせた1つの解」と「付加項のない式の解を初期値の個数だけ組み合わせたもの」の和になる.
\end{bookpoint}

\paragraphbf{形を置いて解く}
答えの形が分かっていれば,未知の定数を含む形を先に置き,条件から定数を決めて解くこともできる.
定数係数の場合,付加項に合わせた1つの解は,まず次の形を候補にできる.
\begin{center}
  \renewcommand{\arraystretch}{1.4}
  \begin{tabular}{l|l}
    付加項&置く形\\ \hline
    定数&定数\\
    $n$の$m$次式&$n$の$m$次式\\
    $q^n$\quad($q\ne0$)&$Aq^n$
  \end{tabular}
\end{center}
この候補が付加項のない式の解と重なるときは,後で見るように$n$を掛けて形を変える.
\sectionref*{節：操作を選ぶ練習}の第1問を,今度は答えの形から見直そう.
例えば
\[
  a_1=1,\qquad a_{n+1}=3a_n+2n+1
\]
なら,付加項のない$b_{n+1}=3b_n$の解は$C\cdot3^{n-1}$である.
付加項は1次式なので,1つの解を$un+v$と置く.
先ほど求めた$u=v=-1$を使うと,\,$-n-1$は元の漸化式を満たす.
したがって$a_n=C\cdot3^{n-1}-n-1$と書け,\,$a_1=C-2=1$より$C=3$なので,
\[
  a_n=3^n-n-1
\]
を得る.\,$a_2=3\cdot1+2+1=6$と$3^2-2-1=6$が一致することも確かめておこう.

\paragraphbf{置いた形がうまくいかないとき}
では,
\[
  a_1=1,\qquad a_{n+1}=2a_n+2^n
\]
はどうだろうか.表にしたがって$A\cdot2^n$と置いてみると,
\[
  A\cdot2^{n+1}=2A\cdot2^n+2^n
\]
となり,\,$0=2^n$という矛盾が出てしまう.
理由は,\,$2^n$がすでに付加項のない式$b_{n+1}=2b_n$の解になっていることにある.
付加項のない式の解をいくら代入しても,付加項は作れないのだ.
こういうときは,\,$n$を1つ掛けた$An\cdot2^{n-1}$と置いてやればよい.代入すると
\[
  A(n+1)2^n=2An\cdot2^{n-1}+2^n
\]
より$A=1$となる.したがって$a_n=C\cdot2^{n-1}+n\cdot2^{n-1}$と書け,\,$a_1=C+1=1$より$C=0$なので,
\[
  a_n=n\cdot2^{n-1}
\]
を得る.両辺を$2^{n+1}$で割って階差型として解いても,同じ答えになる.
このように,付加項の形が付加項のない式の解と重なるときは,\,$n$を掛けた形を試す.
定数も$1^n$と見られるので,特性方程式の解に1がある場合にも同じことが起こる.
例えば$a_{n+1}=a_n+n$では,一次式ではなく二次式$g(n)=\frac{n(n-1)}2$が元の式を満たす.
特性方程式の解が重なっている場合は,重なりの回数だけ$n$を掛ける.\footnote{例えば$a_{n+2}-2a_{n+1}+a_n=1$では,特性方程式の解1は2重であり,\,$g(n)=\frac{n^2}{2}$が元の式を満たす.一般の場合は\sectionref*{節：差分で線形漸化式を解く}を参照.}
その理由は\booknavref{節：付加項の次数が上がる理由}{「付加項の次数が上がる理由」}で詳しく扱う.

\paragraphbf{検算にも使える}
答えの形は,検算にも使える.
例えば付加項のない隣接3項型で,特性方程式が異なる2つの0でない解$\alpha,\beta$を持つときは,答えが$A\alpha^{n-1}+B\beta^{n-1}$の形になる.
重解や付加項がある場合は,それぞれに応じた形を確かめる.
形が合っていても,初期条件と元の漸化式の両方を満たすか検算しよう.

ここで示した答えの形の表は定数係数の場合のものであり,変数係数の場合にはそのまま使えるとは限らない.
一方,2つの解の差を取って付加項を消す操作は,変数係数の線形な漸化式でも使える.
分数型や対数型はそのままでは使えないが,逆数や対数を取って線形な漸化式に直した後なら使える.
本書の解法が全体としてどのような仕組みになっているのかは,付録の最後の\sectionref*{節：本書の解法を見渡す}でまとめる.
